\section{自回归模型（Autoregressive Models）}

\subsection{核心思想：链式分解}

自回归模型利用概率论中的链式法则（Chain Rule）将联合分布分解为一系列条件分布的乘积：

$$p(X) = p(x_1) \cdot p(x_2|x_1) \cdot p(x_3|x_1,x_2) \cdots p(x_T|x_1,\ldots,x_{T-1})$$

\begin{figure}[H]
\centering
\includegraphics[width=0.95\textwidth]{slides-images/slide-019.jpg}
\caption{链式法则：将联合分布 $p(x_1, x_2, \ldots, x_T)$ 分解为条件分布的乘积\protect\footnotemark}
\end{figure}
\footnotetext{Slides 第20页。}

\begin{importantbox}{自回归模型的两个核心要素}
\begin{enumerate}
    \item \textbf{将数据分解为序列}：定义元素之间的顺序（对语言是自然的，对图像需要人为定义）
    \item \textbf{用神经网络建模条件分布}：每一步预测下一个元素的分布 $p(x_t | x_1, \ldots, x_{t-1})$
\end{enumerate}
\end{importantbox}

\subsection{序列建模与采样}

训练时，自回归模型通过 Teacher Forcing 的方式并行计算所有条件概率，然后通过 MLE 优化。采样时，则需要逐元素顺序生成：

\begin{figure}[H]
\centering
\includegraphics[width=0.95\textwidth]{slides-images/slide-021.jpg}
\caption{自回归采样过程：每次生成一个 token，将其追加到输入序列中，然后预测下一个\protect\footnotemark}
\end{figure}
\footnotetext{Slides 第22页。}

\begin{enumerate}
    \item 从 $p(x_1)$ 采样得到 $x_1$
    \item 将 $x_1$ 输入模型，从 $p(x_2|x_1)$ 采样得到 $x_2$
    \item 将 $x_1, x_2$ 输入模型，从 $p(x_3|x_1,x_2)$ 采样得到 $x_3$
    \item 重复直到序列完成
\end{enumerate}

\subsection{语言建模}

自回归模型天然适合语言建模，因为语言本身就是一维的离散序列。

\begin{figure}[H]
\centering
\includegraphics[width=0.95\textwidth]{slides-images/slide-024.jpg}
\caption{语言建模中的自回归：用 RNN 或 Transformer 建模 token 序列的条件概率\protect\footnotemark}
\end{figure}
\footnotetext{Slides 第25页。}

序列中的每个元素是一个 token（可以是字符、子词或单词），模型输出每个位置上所有可能 token 的概率分布（通常通过 softmax）。ChatGPT 等大语言模型本质上就是超大规模的自回归模型。

\begin{figure}[H]
\centering
\includegraphics[width=0.95\textwidth]{slides-images/slide-026.jpg}
\caption{温度采样：温度参数 $\tau$ 控制采样的随机性。$\tau \to 0$ 趋向贪心解码，$\tau \to \infty$ 趋向均匀随机\protect\footnotemark}
\end{figure}
\footnotetext{Slides 第27页。}

$$p(x_t = v | x_{<t}) = \frac{\exp(s_v / \tau)}{\sum_k \exp(s_k / \tau)}$$

\begin{itemize}
    \item $s_v$：模型输出的 logit
    \item $\tau$：温度参数
    \item $\tau \to 0$：趋向 argmax（确定性，低多样性）
    \item $\tau \to \infty$：趋向均匀分布（高多样性，低质量）
\end{itemize}

\subsection{图像上的自回归模型}

将图像展平为像素序列，每个子像素（R/G/B）是 0-255 的离散整数值。一张 $1024 \times 1024$ 的图片会产生约 300 万个 token 的序列，计算代价极高。

\begin{figure}[H]
\centering
\includegraphics[width=0.95\textwidth]{slides-images/slide-029.jpg}
\caption{PixelRNN/PixelCNN：将图像像素视为序列，逐像素进行自回归建模\protect\footnotemark}
\end{figure}
\footnotetext{Slides 第30页。}

\begin{warningbox}{直接在像素空间做自回归效率极低}
一张 $1024 \times 1024$ 的 RGB 图像有 $\sim$3M 子像素值。这意味着自回归序列的长度达到 300 万，对 Transformer 而言计算代价极其昂贵。解决方案是先将图像编码为更短的 latent token 序列，然后在 latent 空间上做自回归——这正是现代 token-based 图像生成方法的核心思路。
\end{warningbox}

\begin{figure}[H]
\centering
\includegraphics[width=0.95\textwidth]{slides-images/slide-032.jpg}
\caption{PixelCNN 生成的图像样本：分辨率低，质量有限，但展示了自回归方法的可行性\protect\footnotemark}
\end{figure}
\footnotetext{Slides 第33页。}

\subsection{本章小结}

自回归模型利用链式法则将联合分布分解为条件分布的乘积，可以精确计算 $p(X)$ 并用 MLE 训练。它天然适合语言建模，但在图像上面临序列过长的问题。现代方法通过 latent tokenization 缓解这一问题。

%% ========== Part 5: 变分自编码器 ==========

